\documentclass [11pt,twoside]{article}
\usepackage[utf8]{inputenc}
\usepackage[T1]{fontenc}

%Page margins, header and footer positions
\usepackage{geometry}
 \geometry{
 a4paper,
 total={210mm,297mm},
 left=25mm,
 right=25mm,
 top=30mm,
 bottom=25mm,
 headsep=7mm}

\interfootnotelinepenalty=10000

%To display filling dots in the TOC for all entries
\usepackage[titles]{tocloft}
\renewcommand{\cftsecleader}{\cftdotfill{\cftdotsep}}

%Define new header and footer style
\usepackage{fancyhdr}

\pagestyle{fancy}
\fancyhf{}
\lhead{\color{Gray}{\small{}}}
\lfoot{\textcolor{Gray}{\small{}}}
\rfoot{\textcolor{Gray}{\thepage}}
\renewcommand{\headrulewidth}{0pt}

%PACKAGES
\usepackage{wasysym}
\usepackage{pifont}

\newcommand{\supported}{\ding{52}\xspace}
\newcommand{\unsupported}{\ding{55}\xspace}
\newcommand{\partsupported}{\textcolor{black!40}{\ding{52}}\xspace}
\newcommand{\lowsupported}{\textcolor{black!20}{\ding{52}}\xspace}
\newcommand{\unknowsupported}{\textbf{?}\xspace}

%Font: Times
\usepackage{times}
%Change monospaced font
\renewcommand{\ttdefault}{lmtt}

%tables
\usepackage{tabu}
\usepackage{tabularx}
\usepackage{ltablex}
\usepackage{longtable}
\usepackage{float} % To allow the use of H modifier in long tables

%landscape mode
\usepackage{pdflscape}
\usepackage{rotating}
\usepackage{caption}

%make landscape mode be sensitive to even and odd pages
%start
\def\myrotate{\ifodd\c@page\else-\fi 90}
\makeatletter
\global\let\orig@begin@landscape=\landscape%
\global\let\orig@end@landscape=\endlandscape%
\gdef\@true{1}
\gdef\@false{0}
\gdef\landscape{%
    \global\let\within@landscape=\@true%
    \orig@begin@landscape%
}%
\gdef\endlandscape{%
    \orig@end@landscape%
    \global\let\within@landscape=\@false%
}%
\@ifpackageloaded{pdflscape}{%
    \gdef\pdf@landscape@rotate{\PLS@Rotate}%
}{
    \gdef\pdf@landscape@rotate#1{}%
}
\let\latex@outputpage\@outputpage
\def\@outputpage{
    \ifx\within@landscape\@true%
        \if@twoside%
            \ifodd\c@page%
                \gdef\LS@rot{\setbox\@outputbox\vbox{%
                    \pdf@landscape@rotate{-90}%
                    \hbox{\rotatebox{90}{\hbox{\rotatebox{180}{\box\@outputbox}}}}}%
                }%
            \else%
                \gdef\LS@rot{\setbox\@outputbox\vbox{%
                    \pdf@landscape@rotate{+90}%
                    \hbox{\rotatebox{90}{\hbox{\rotatebox{0}{\box\@outputbox}}}}}%
                }%
            \fi%
        \else%
            \gdef\LS@rot{\setbox\@outputbox\vbox{%
                \pdf@landscape@rotate{+90}%
                \hbox{\rotatebox{90}{\hbox{\rotatebox{0}{\box\@outputbox}}}}}%
            }%
        \fi%
    \fi%
    \latex@outputpage%
}
\makeatother
%end

%graphics
\usepackage{graphicx}
\usepackage[dvipsnames, table]{xcolor}
%If you upload images from PC, you need to insert code for the path here (different for Windows and Unix OS)

%References
%\usepackage{xpatch}
%\usepackage[backend=biber, style=numeric, citestyle=numeric, sorting=none]{biblatex}
%\addbibresource{main.bib}

%Other
\usepackage{ifthen}
\usepackage{xspace}
\usepackage{enumitem}
\usepackage{amssymb}
\usepackage[pdftex, colorlinks]{hyperref}
\newcommand{\comment}[1]{{\color{Red}$\blacktriangleright$ Comment: #1 $\blacktriangleleft$}}


% Some utilities\ldots
\usepackage{soul}
\usepackage{tikz}

\usetikzlibrary{calc}
\usetikzlibrary{decorations.pathmorphing}


\makeatletter

\newcommand{\defhighlighter}[3][]{%
  \tikzset{every highlighter/.style={color=#2, fill opacity=#3, #1}}%
}

\defhighlighter{yellow}{.5}

\newcommand{\highlight@DoHighlight}{
  \fill [ decoration = {random steps, amplitude=1pt, segment length=15pt}
        , outer sep = -15pt, inner sep = 0pt, decorate
       , every highlighter, this highlighter ]
        ($(begin highlight)+(0,8pt)$) rectangle ($(end highlight)+(0,-3pt)$) ;
}

\newcommand{\highlight@BeginHighlight}{
  \coordinate (begin highlight) at (0,0) ;
}

\newcommand{\highlight@EndHighlight}{
  \coordinate (end highlight) at (0,0) ;
}

\newdimen\highlight@previous
\newdimen\highlight@current

\DeclareRobustCommand*\highlight[1][]{%
  \tikzset{this highlighter/.style={#1}}%
  \SOUL@setup
  %
  \def\SOUL@preamble{%
    \begin{tikzpicture}[overlay, remember picture]
      \highlight@BeginHighlight
      \highlight@EndHighlight
    \end{tikzpicture}%
  }%
  %
  \def\SOUL@postamble{%
    \begin{tikzpicture}[overlay, remember picture]
      \highlight@EndHighlight
      \highlight@DoHighlight
    \end{tikzpicture}%
  }%
  %
  \def\SOUL@everyhyphen{%
    \discretionary{%
      \SOUL@setkern\SOUL@hyphkern
      \SOUL@sethyphenchar
      \tikz[overlay, remember picture] \highlight@EndHighlight ;%
    }{%
    }{%
      \SOUL@setkern\SOUL@charkern
    }%
  }%
  %
  \def\SOUL@everyexhyphen##1{%
    \SOUL@setkern\SOUL@hyphkern
    \hbox{##1}%
    \discretionary{%
      \tikz[overlay, remember picture] \highlight@EndHighlight ;%
    }{%
    }{%
      \SOUL@setkern\SOUL@charkern
    }%
  }%
  %
  \def\SOUL@everysyllable{%
    \begin{tikzpicture}[overlay, remember picture]
      \path let \p0 = (begin highlight), \p1 = (0,0) in \pgfextra
        \global\highlight@previous=\y0
        \global\highlight@current =\y1
      \endpgfextra (0,0) ;
      \ifdim\highlight@current < \highlight@previous
        \highlight@DoHighlight
        \highlight@BeginHighlight
      \fi
    \end{tikzpicture}%
    \the\SOUL@syllable
    \tikz[overlay, remember picture] \highlight@EndHighlight ;%
  }%
  \SOUL@
}

\makeatother

% Common abbrev. are set as commands to ensure proper spacing after the dot
\RequirePackage{xspace}
\newcommand{\ie}{i.e.\@\xspace}
\newcommand{\aka}{a.k.a.\@\xspace}
\newcommand{\Ie}{I.e.\@\xspace}
\newcommand{\cf}{cf.\@\xspace}
\newcommand{\Cf}{Cf.\@\xspace}
\newcommand{\eg}{e.g.\@\xspace}
\newcommand{\Eg}{E.g.\@\xspace}
\newcommand{\etal}{et al.\@\xspace}
\newcommand{\etc}{etc.\@\xspace}
\newcommand{\wrt}{w.r.t.\@\xspace}
\newcommand{\Wrt}{W.r.t.\@\xspace}



\date{}


\usepackage{graphicx}
\usepackage{float}
\usepackage{listings}
\usepackage{xcolor}
\usepackage[textsize=footnotesize]{todonotes}

\hypersetup{
    colorlinks=true,
    linkcolor=black,
    urlcolor=black,
    citecolor=black
}

\graphicspath{{Images/}}
\begin{document}
\setlength{\parindent}{0pt}
\setlength{\parskip}{6pt}

%TITLE PAGE

\begin{titlepage}

{\begin{table}[t!]
\centering
\begin{tabu} to \textwidth { X[1.3,r,p] X[1.7,l,p] }
\textcolor{Blue}
{\textbf{\small{Hoxha Eldi, Prunescu Bogdan-Andrei, Marian Sorin-Nicu}}} & \includegraphics[scale=0.5]{Images/PolimiLogo}
\end{tabu}
\end{table}}~\\ [7cm]

\begin{flushleft}

{\textcolor{Blue}{\textbf{\Huge{TAVI: Trajectory-Level Verification of Embodied AI Agents via LTL}}}} \\ [1cm]

\end{flushleft}

\end{titlepage}

\begin{table}[h!]
\begin{tabu} to \textwidth { X[0.3,r,p] X[0.7,l,p] }
\hline

\textbf{Deliverable:} & Brief Scientific Report \\
\textbf{Title:} & TAVI: Trajectory-Level Verification of Embodied AI Agents via LTL \\
\textbf{Authors:} & Hoxha Eldi, Prunescu Bogdan-Andrei, Marian Sorin-Nicu
 \\
\textbf{Supervisor:} & Matteo Camilli \\
\textbf{Version:} & 1 \\
\textbf{Date:} & % TODO
\\
\hline
\end{tabu}
\end{table}

\setcounter{page}{2}

%------------------------------------------------------------------------------------------------------------------------------------------------
\newpage
\addcontentsline{toc}{section}{Table of Contents}
\tableofcontents
\newpage
\addcontentsline{toc}{section}{List of Figures}
\listoffigures
\addcontentsline{toc}{section}{List of Tables}
\listoftables
\newpage

%------------------------------------------------------------------------------------------------------------------------------------------------
{\color{Blue}{\section{Introduction}}}
\todo[inline]{Motivate the problem: per-step checking (AgentSpec) misses hazards that only show
up across a trajectory (e.g. an object made unsafe by an earlier action). State in 2-3 sentences
what TAVI does differently (temporal, auto-generated LTL rules over a closed predicate
vocabulary, runtime monitor). End with a one-sentence roadmap of the rest of the report.}

%------------------------------------------------------------------------------------------------------------------------------------------------
{\color{Blue}{\section{Related Work}}}
\todo[inline]{One short paragraph each on: (1) AgentSpec -- trigger/check/enforce DSL, per-step
only, manually authored rules; (2) SafeAgentBench -- safety-aware task planning benchmark,
evaluates agent behavior, not a checking/monitoring mechanism itself, this is our evaluation
benchmark (Section 4); (3) Embodied AI Agents (Fung et al.) -- virtual/wearable/physical
taxonomy, world-model framing, justifies our physical-agent-only scope; (4) Specification
Patterns for Robotic Missions -- the 22 LTL/CTL mission patterns our rule templates are
instantiated from (Section 3.3). Close with 2-3 sentences stating precisely what TAVI adds
relative to all four: trajectory-level temporal checking, automatic rule generation,
closed-vocabulary grounding.}

%------------------------------------------------------------------------------------------------------------------------------------------------
{\color{Blue}{\section{Method}}}

{\color{Blue}{\subsection{Pipeline Overview}}}
\todo[inline]{Describe the 5 conceptual modules (trace collector, rule generation, trace
translation/grounding, LTL checking, runtime monitor); state which are implemented end-to-end
vs. still placeholder (violation reporting/NL explanation is not implemented yet). Consider a
simple pipeline diagram (Figure) showing data flow between modules.}

{\color{Blue}{\subsection{Trace Collection and Predicate Vocabulary}}}
\todo[inline]{Describe the JSONL trace schema (step id, timestamp, action, pre/post state,
observation, propositions, llm\_call, tool\_calls, outcome). Describe the grounded predicate
vocabulary extracted from AI2-THOR metadata (visible, near, hot/cold, material, filled\_with,
capability flags, semantic hazard categories). Explain the closed-vocabulary contract
(trace\_predicates.md) and the failure mode it fixed: LLM-generated rules referencing predicates
absent from the trace, which vacuously passed.}

{\color{Blue}{\subsection{Rule Generation}}}
\todo[inline]{Describe online, per-task, cached LTL rule generation via Gemini 2.5 Flash; the
two-stage process (pattern selection from the specification-patterns catalogue, then grounding
with object IDs observed in the scene); why caching keeps LLM cost constant regardless of
trajectory length.}

{\color{Blue}{\subsection{LTL Checking}}}
\todo[inline]{Describe flloat/LTLf checking mechanics, what ``violation'' means operationally
over a finite trace, and how the runtime monitor (main.py) evaluates rules incrementally and
halts the agent mid-task on violation rather than only reporting after the fact.}

{\color{Blue}{\subsection{Implementation Notes}}}
\todo[inline]{Pick 2-3 concrete engineering details worth reporting: the agent\_executed
detection fix (planner error strings no longer counted as executed actions); the parallel
benchmark runner (run\_phases.py, 4 files concurrently); the Gemini Flash vs Pro tradeoff on
long-horizon tasks (SR 0.86 vs 0.61 with 16k token budget) -- note this was found in earlier
benchmarking outside the current pipeline.}

%------------------------------------------------------------------------------------------------------------------------------------------------
{\color{Blue}{\section{Evaluation Setup}}}
\todo[inline]{Describe SafeAgentBench (750 tasks, 10 hazard types), the 4 task categories used
(unsafe/safe/abstract/long-horizon) with 1-line definitions, sample size (100 tasks, 25 per
category), and the outcome taxonomy (TP/TN/FN/FP/Agent Avoided Action/No Trace) with 1-line
definitions of each as a table. Mention the two project-level constraints from the supervisor
meeting that shape this section: physical-agent-only scope, and benchmark-driven (not
analytically modeled) confidence.}

%------------------------------------------------------------------------------------------------------------------------------------------------
{\color{Blue}{\section{Results}}}
\todo[inline]{Insert the 100-task results table (TP 23/39, TN 50/50, FN 16, FP 0, Agent Avoided
Action 2, No Trace 9) and headline metrics (82\% accuracy, 59\% recall, 100\% precision).
Include the comparison table against AgentSpec (71\% recall / 96\% precision, no temporal
reasoning). Also report the earlier 40-task run (92.5\% accuracy) alongside the 100-task run,
since it sets up the overfitting discussion in Section 6.}

%------------------------------------------------------------------------------------------------------------------------------------------------
{\color{Blue}{\section{Discussion and Limitations}}}
\todo[inline]{Cover: overfitting as task diversity increases (92.5\% on 40 tasks to 82\% on 100
tasks); the recall bottleneck of 16 missed unsafe tasks despite 100\% precision, worst on the
abstract category where the hazard is described indirectly; the 9 no-trace tasks being
agent/planner crashes upstream of the checker (not checker failures) that still reduce effective
coverage; confidence intervals must be benchmark-driven per supervisor constraint and are not
yet exposed (pipeline currently reports a boolean satisfied/violated verdict only); and the scope
being limited to physical embodied agents per project constraints.}

%------------------------------------------------------------------------------------------------------------------------------------------------
{\color{Blue}{\section{Conclusion and Next Steps}}}
\todo[inline]{One paragraph summarizing the contribution and headline results, then concrete
next steps: diagnose specific abstract/long-horizon failure cases (vocabulary too narrow vs. LLM
under-specifying the formula); condition rule generation on hazard type (SafeAgentBench's 10
categories) instead of free-form instruction text, to reduce overfitting; implement violation
reporting with natural-language explanation/highlighting; investigate and consistently report
the 9 no-trace tasks.}

%------------------------------------------------------------------------------------------------------------------------------------------------
{\color{Blue}{\section{References}}}

\todo[inline]{Papers in Docs/Papers/ -- cite these first. Fill in full bibliographic details
(authors, venue, year, pages) for each before submission.}
\begin{enumerate}
    \item AgentSpec: Customizable Runtime Enforcement for Safe and Reliable LLM Agents. \todo{authors, venue, year}
    \item SafeAgentBench: a benchmark for safety-aware task planning of embodied LLM agents. \todo{authors, venue, year}
    \item Embodied AI Agents: Modeling the World. \todo{authors (Fung et al.), venue, year}
    \item Specification Patterns for Robotic Missions. \todo{authors, venue, year}
\end{enumerate}

\todo[inline]{Other sources researched during the project (tools, libraries, frameworks) -- add
as used, with full citation/URL for each we actually relied on.}
\begin{itemize}
    \item AI2-THOR simulation environment. \todo{paper/URL}
    \item LangGraph agent framework. \todo{docs/URL}
    \item flloat LTLf checking library. \todo{paper/URL}
    \item ProbGuard -- LangChain-based probabilistic runtime monitoring. \todo{mentioned
    2026-04-01 meeting; confirm if actually used or drop}
    \item Gemini 2.5 Flash / Pro model documentation. \todo{URL}
\end{itemize}

\listoftodos

\end{document}
